%! Composition of Functions — Unit 5, Lesson 5.6 — Group Activity (Answer Key)
\documentclass[10pt]{article}
\usepackage{algebra2-article}
\usepackage{algebra2-key}   % pulls in -boxes; do NOT also load -boxes
\newcommand{\TallMath}[1]{$\displaystyle #1\rule[-1.4em]{0pt}{3.2em}$}
% In-formula answer slots (\ans is text-mode and must never appear inside $...$).
\newcommand{\mans}[1]{{\color{keyred}\mathbf{#1}}}
\newcommand{\mexp}[1]{{\color{keyred}#1}}
\newcommand{\lgrid}[4]{%
  \draw[step=1, linegray!40, very thin] (#1,#3) grid (#2,#4);
  \draw[->, semithick, charcoal] (#1,0)--(#2,0) node[below right, font=\scriptsize]{$x$};
  \draw[->, semithick, charcoal] (0,#3)--(0,#4) node[left, font=\scriptsize]{$y$};}

\begin{document}

\pageheader{Unit 5, Lesson 5.6}{Group Activity --- Answer Key}
\namepartnerperiod

\vspace{0.08in}

\begin{headlinebox}{forestbg}
\small\textbf{Wired in Series.} Two machines, one after the other --- and the wire between them
carries exactly one number. With your partner, write that middle number down \emph{every single
time}. Most mistakes today are not arithmetic; they are running the machines in the wrong order.
\end{headlinebox}

\vspace{0.06in}

\begin{tcolorbox}[colback=white, colframe=black!40, title=\textbf{Tier R --- Remediate},
    fonttitle=\bfseries, arc=1.5mm, left=3mm, right=3mm, top=2mm, bottom=2mm, breakable]
\textbf{Part 1 --- One step at a time.} \; Let $f(x)=x+4$ and $g(x)=3x$.
\begin{center}
\renewcommand{\arraystretch}{1.8}
\begin{tabularx}{0.96\linewidth}{@{}l X X X@{}}
\hline
 & \textbf{Which runs first?} & \textbf{The middle number} & \textbf{Final answer} \\
\hline
$(f\circ g)(2)$ & \ans{$g$} & $g(2)=\mans{6}$ & $f(\mans{6})=\mans{10}$ \\
$(g\circ f)(2)$ & \ans{$f$} & $f(2)=\mans{6}$ & $g(\mans{6})=\mans{18}$ \\
\hline
\end{tabularx}
\end{center}
\vspace{-4pt}
Same two functions, same input, two answers. Which one is bigger, and why does the order do that?
\ans{$(g\circ f)(2)=18$. When tripling happens \emph{last}, the $+4$ gets tripled too.}

\vspace{6pt}
\textbf{Part 2 --- Now you run them.} Let $f(x)=2x-1$ and $g(x)=x^{2}$. Show the middle number.
\begin{center}
\renewcommand{\arraystretch}{1.8}
\begin{tabularx}{\linewidth}{@{}X X@{}}
(a) \; $(f\circ g)(3)$: \; middle $=\mans{9}$, \; answer $=\mans{17}$ &
(b) \; $(g\circ f)(3)$: \; middle $=\mans{5}$, \; answer $=\mans{25}$ \\
(c) \; $(f\circ g)(-2)$: \; middle $=\mans{4}$, \; answer $=\mans{7}$ &
(d) \; $(g\circ f)(0)$: \; middle $=\mans{-1}$, \; answer $=\mans{1}$ \\
\end{tabularx}
\end{center}

\vspace{4pt}
\textbf{Part 3 --- Read it off the graphs.} These are the same two graphs from your notes.
\begin{center}
\begin{minipage}[c]{0.47\linewidth}
\centering
\begin{tikzpicture}[scale=0.34]
  \lgrid{-5}{7}{-3}{5}
  \draw[very thick, forest] (-4,-2)--(0,2)--(4,0)--(6,4);
  \fill[forest] (0,2) circle (3.4pt);
  \fill[forest] (4,0) circle (3.4pt);
  \node[forest, font=\scriptsize] at (-2.6,2.6) {$y=f(x)$};
\end{tikzpicture}
\end{minipage}
\hfill
\begin{minipage}[c]{0.47\linewidth}
\centering
\begin{tikzpicture}[scale=0.34]
  \lgrid{-5}{7}{-3}{5}
  \draw[very thick, navy] (-4,4)--(2,-2)--(6,2);
  \fill[navy] (2,-2) circle (3.4pt);
  \node[navy, font=\scriptsize] at (4.4,3.4) {$y=g(x)$};
\end{tikzpicture}
\end{minipage}
\end{center}
\vspace{-2pt}
\begin{center}
\renewcommand{\arraystretch}{1.7}
\begin{tabularx}{0.96\linewidth}{@{}X X X X@{}}
\hline
 & \textbf{Middle number} & \textbf{Read on which graph?} & \textbf{Answer} \\
\hline
(a) $(f\circ g)(-4)$ & \ans{$4$} & \ans{$g$} & \ans{$0$} \\
(b) $(g\circ f)(-4)$ & \ans{$-2$} & \ans{$f$} & \ans{$2$} \\
(c) $(f\circ g)(6)$ & \ans{$2$} & \ans{$g$} & \ans{$1$} \\
(d) $(g\circ f)(2)$ & \ans{$1$} & \ans{$f$} & \ans{$-1$} \\
\hline
\end{tabularx}
\end{center}
\vspace{-4pt}
Rows (a) and (b) start from the same input, $x=-4$. Explain to your partner why they do not finish
at the same place. \ans{They hand different middle numbers along the wire: (a) sends $g(-4)=4$ into
$f$, (b) sends $f(-4)=-2$ into $g$.}
\end{tcolorbox}

\begin{tcolorbox}[colback=white, colframe=black!40, title=\textbf{Tier A --- Approaching Proficiency},
    fonttitle=\bfseries, arc=1.5mm, left=3mm, right=3mm, top=2mm, bottom=2mm, breakable]
\textbf{Part 1 --- Both orders, in general.} For each pair, find both compositions and then decide
whether they are the same function.
\begin{center}
\renewcommand{\arraystretch}{1.8}
\begin{tabularx}{\linewidth}{@{}l X X c@{}}
\hline
\textbf{$f$ and $g$} & \textbf{$(f\circ g)(x)$} & \textbf{$(g\circ f)(x)$} & \textbf{Same?} \\
\hline
(a) $f(x)=3x$, \; $g(x)=x-2$ & \ans{$3x-6$} & \ans{$3x-2$} & \ans{no} \\
(b) $f(x)=x^{2}$, \; $g(x)=x+1$ & \ans{$x^{2}+2x+1$} & \ans{$x^{2}+1$} & \ans{no} \\
(c) $f(x)=\sqrt{x}$, \; $g(x)=x-9$ & \ans{$\sqrt{x-9}$} & \ans{$\sqrt{x}-9$} & \ans{no} \\
(d) $f(x)=x+3$, \; $g(x)=x-3$ & \ans{$x$} & \ans{$x$} & \ans{yes} \\
(e) $f(x)=3x-1$, \; $g(x)=\dfrac{x+1}{3}$ & \ans{$x$} & \ans{$x$} & \ans{yes} \\
\hline
\end{tabularx}
\end{center}
\vspace{1.35in}

\textbf{Part 2 --- What did you just find?}
\begin{enumerate}[label=\textbf{J\arabic*.}, leftmargin=2.4em, itemsep=6pt]
  \item Rows (d) and (e) behaved differently from the others: \emph{both} orders came out to the
        same thing, and that thing was remarkably plain. What was it? What are those two functions
        doing to each other --- and what word would you invent for a pair like that?
        \ansline{Both compositions came out to $x$ itself --- whatever you start with, you get back.
        Each function \emph{undoes} the other: $g$ subtracts $3$ and $f$ adds it back; $g$ divides
        the shifted number by $3$ and $f$ multiplies and shifts back. Student wording will vary
        (``opposites,'' ``reversers,'' ``undo-ers''); the real word is \textbf{inverse functions},
        which is Lesson 5.7. Note for them that the check $f(g(x))=g(f(x))=x$ is exactly what they
        just performed.}
  \item Row (c) is the one to be careful with. The two answers do not have the same domain. State
        both domains, and for each one say \emph{which stage} of the composition does the refusing.
        \ansline{$\sqrt{x-9}$ has domain $x\ge9$: stage 1 accepts everything, but stage 2 (the
        square root) demands $x-9\ge0$. \; $\sqrt{x}-9$ has domain $x\ge0$: here the square root
        \emph{is} stage 1, so it refuses negatives immediately; subtracting $9$ afterwards restricts
        nothing.}
  \item A classmate answered row (b) with $(f\circ g)(x)=x^{2}+1$. That is a real function --- just
        not the one that was asked for. Which composition did they actually compute, and what is the
        one-sentence fix? \ansline{They computed $(g\circ f)(x)$ --- squared first, added $1$
        second. Fix: the function next to the $x$ runs first, so $f$ receives the \emph{whole}
        expression $x+1$ and squares all of it: $(x+1)^{2}$.}
\end{enumerate}

\vspace{4pt}
\textbf{Part 3 --- Domains, on purpose.} Find $(f\circ g)(x)$, simplify it, and \emph{then} work out
the domain from the two machines.
\begin{center}
\renewcommand{\arraystretch}{1.8}
\begin{tabularx}{\linewidth}{@{}l X X@{}}
\hline
\textbf{$f$ and $g$} & \textbf{$(f\circ g)(x)$, simplified} & \textbf{Domain} \\
\hline
(a) $f(x)=\sqrt{x}$, \; $g(x)=x+2$ & \ans{$\sqrt{x+2}$} & \ans{$x\ge-2$} \\
(b) $f(x)=\sqrt{x-1}$, \; $g(x)=4x$ & \ans{$\sqrt{4x-1}$} & \ans{$x\ge\tfrac14$} \\
(c) $f(x)=x^{2}+1$, \; $g(x)=\sqrt{x}$ & \ans{$x+1$} & \ans{$x\ge0$} \\
\hline
\end{tabularx}
\end{center}
\vspace{-2pt}
Row (c) simplifies to something with no radical left in it at all. Why is its domain still
restricted? \ansline{Because the restriction happens at \emph{Gate 1}: the inner machine is
$\sqrt{x}$, and it refuses a negative input before the squaring ever runs. Squaring afterwards
cleans up the radical in the \emph{formula}, but it cannot go back and admit inputs that were
already turned away.}
\end{tcolorbox}

\begin{tcolorbox}[colback=white, colframe=black!40, title=\textbf{Tier E --- Extension},
    fonttitle=\bfseries, arc=1.5mm, left=3mm, right=3mm, top=2mm, bottom=2mm, breakable]
\textbf{Part 1 --- Settle the jacket argument for every price.} Keep $f(x)=0.80x$ (the $20\%$
discount) and $g(x)=x-10$ (the coupon), but stop using \$$80$.
\begin{enumerate}[label=(\alph*), leftmargin=1.9em, itemsep=6pt]
  \item Find both compositions as rules. \; $(f\circ g)(x)=\ans{$0.8x-8$}$ \; and \;
        $(g\circ f)(x)=\ans{$0.8x-10$}$
        \vspace{0.35in}
  \item Subtract one from the other. The gap between the two orders is \ans{\$$2$} --- and it does
        not depend on the price at all. Explain where that number comes from in terms of the
        \$$10$ coupon and the $20\%$ discount. \ansline{Applying the coupon first means the $20\%$
        discount is also taken off the \$$10$ you already saved: $0.20\times10=\$2$ of the coupon is
        discounted away. The $x$-terms cancel in the subtraction, so the loss is \$$2$ at every
        price.}
  \item Which order does the \emph{store} prefer, and which does the \emph{customer} prefer? Does
        your answer change for an expensive jacket? \ansline{The store prefers coupon-first (the
        customer pays \$$2$ more); the customer prefers discount-first. It never changes --- the gap
        is the constant $2$, not a percentage of the price.}
\end{enumerate}

\vspace{4pt}
\textbf{Part 2 --- When \emph{do} two linear functions commute?} Let $f(x)=ax+b$ and $g(x)=cx+d$.
\begin{enumerate}[label=(\alph*), leftmargin=1.9em, itemsep=6pt, start=4]
  \item Compose both ways and simplify. \; $(f\circ g)(x)=\ans{$acx+ad+b$}$, \;
        $(g\circ f)(x)=\ans{$acx+cb+d$}$
        \vspace{0.40in}
  \item The $x$-terms are identical no matter what. So the two compositions are equal for
        \emph{every} $x$ exactly when the constants match: \ans{$ad+b=cb+d$}
  \item Test your condition on Tier A row (d) ($f(x)=x+3$, $g(x)=x-3$) and on the jacket pair
        ($f(x)=0.8x$, $g(x)=x-10$). Does it predict the right verdict both times? \ansline{Row (d):
        $a=1,b=3,c=1,d=-3$, so $ad+b=0$ and $cb+d=0$ --- equal, and indeed both compositions gave
        $x$. Jacket: $a=0.8,b=0,c=1,d=-10$, so $ad+b=-8$ and $cb+d=-10$ --- not equal, and the gap
        between them, $-8-(-10)=2$, is the \$$2$ from Part 1.}
\end{enumerate}

\vspace{4pt}
\textbf{Part 3 --- More than two machines.}
\begin{enumerate}[label=(\alph*), leftmargin=1.9em, itemsep=6pt, start=7]
  \item Let $f(x)=x^{2}$, $g(x)=x-1$, $h(x)=3x$. Find $(f\circ g\circ h)(2)$ by going right to left.
        \; middles: \ans{$6$}, \ans{$5$} \; answer $=\ans{$25$}$
  \item Take $h(x)=2\sqrt{x-5}+1$ apart into \emph{three} machines run in order. \\
        First: \ans{$x-5$} \quad Second: \ans{$\sqrt{x}$} \quad Third: \ans{$2x+1$}
  \item In Lesson 5.4 you called $2\sqrt{x-5}+1$ ``the parent, shifted right $5$, stretched by $2$,
        shifted up $1$.'' Compare that list with your answer to (h). What is a transformation,
        really? \ansline{It is the same list. A transformation \emph{is} a composition: the inside
        machine changes the input before the parent sees it (horizontal moves), and the outside
        machine changes the parent's answer afterwards (vertical moves).}
\end{enumerate}

\vspace{4pt}
\textbf{Part 4 --- The pair that almost undoes itself.} Complete both columns. Write ``undefined''
where the expression does not exist over the real numbers.
\begin{center}
\renewcommand{\arraystretch}{1.6}
\begin{tabularx}{0.9\linewidth}{@{}c X X@{}}
\hline
\textbf{$x$} & \TallMath{\left(\sqrt{x}\,\right)^{2}} & \TallMath{\sqrt{x^{2}}} \\
\hline
$-4$ & \ans{undefined} & \ans{$4$} \\
$-1$ & \ans{undefined} & \ans{$1$} \\
$0$ & \ans{$0$} & \ans{$0$} \\
$4$ & \ans{$4$} & \ans{$4$} \\
$9$ & \ans{$9$} & \ans{$9$} \\
\hline
\end{tabularx}
\end{center}
\begin{enumerate}[label=(\alph*), leftmargin=1.9em, itemsep=5pt, start=10]
  \item As formulas, $\left(\sqrt{x}\,\right)^{2}=\ans{$x$}$ and $\sqrt{x^{2}}=\ans{$|x|$}$
        (Lesson 5.1). One of them is defined for every real number; the other is not.
        \ans{$\sqrt{x^{2}}$ is defined everywhere; $\left(\sqrt{x}\,\right)^{2}$ needs $x\ge0$.}
  \item On which inputs do \emph{both} columns hand you back the $x$ you started with?
        \ans{$x\ge0$}
  \item So squaring and square-rooting undo each other --- but only on part of the number line.
        Tomorrow you will meet the word for two functions that undo each other everywhere. What do
        you predict has to be done to $y=x^{2}$ first, to make that possible? \ansline{Its domain
        has to be \emph{restricted} to $x\ge0$ (keep only the right half of the parabola). On that
        half, squaring and square-rooting really are perfect opposites in both orders --- which is
        the answer to the question left open in Lesson 5.0.}
\end{enumerate}
\end{tcolorbox}

\begin{teachernote}
\small \textbf{Assign one tier, not all three.} Tier R is a procedure clinic; Tier A is the target
work and every student headed for the exit ticket should get through its Part 1 and J2; Tier E is a
genuine extension and its Part 4 is the direct set-up for Lesson 5.7.
\begin{itemize}[leftmargin=1.3em, nosep]
  \item \textbf{Insist on the middle number} in Tier R. A pair that stops writing it is a pair about
        to reverse an order. Tier R Part 3 rows (a) and (b) exist only to make that visible: same
        starting input, different wire, different destination.
  \item \textbf{Tier A rows (d) and (e) are the point of the whole activity.} Do \emph{not} tell
        groups what they have found. Let J1 be answered in their own words and collect the wordings
        --- ``they cancel,'' ``they undo each other,'' ``you get $x$ back'' --- and open Lesson 5.7
        by putting those exact phrases on the board next to the word \emph{inverse}. Row (e) matters
        because it is the first pair students could not have guessed by eye.
  \item \textbf{J2 is the item to grade closely.} ``Different domains'' is not the answer; naming
        \emph{which stage} refuses, and \emph{what number} it refuses, is. Students who cannot do
        this will not follow 5.7's restricted-domain argument.
  \item \textbf{Tier A Part 3(c) is the sleeper.} $x+1$ with domain $x\ge0$ looks like a mistake to
        students. It is the same phenomenon as Tier E Part 4 in smaller print --- if a group finds
        it and argues about it, that argument is worth more than finishing the page.
  \item \textbf{Tier E Part 2} is the only fully general result in the lesson, and it is worth
        pointing out that the two $x$-terms come out identical \emph{automatically}: composition of
        linear functions always multiplies the slopes, which is why only the constants can disagree.
\end{itemize}
\end{teachernote}

\end{document}
